\cotefile{134-1}
\scanpage{2}

Lettres Breen 75, 76\unsure{.}

% archive sheet number: 1

\scanpage{4}

\noindent\hfill Lodève le 27.\unsure{8.74|8/74|8/4}\\
Cher Deligne,\\
% margin note written sideways
\margin{Cher Breen, J'imagine que tu sais ce que le \unsure{maître|maitre} demande, mais n'ai pu \unsure{neamoins|neanmoins} lui \unsure{repondre|répondre} --\\
Qu'en \unsure{pens|penses}-tu. Bien à toi\\
P. \unsure{Deligne}\\}
Étant peut-être empêché \unsure{par ma|par une|par mon} \unsure{faculté|famille|absence} d'\unsure{assurer} un cours de 1\textsuperscript{er} cycle au 1\textsuperscript{er} trimestre, je vais peut\unsure{-|'-}être à la place faire un petit séminaire d'algèbre, et envisage de le faire sur les \unsure{fourbis} de Mme \unsure{Sinh}, éventuellement transposés dans le contexte des ``champs''. À ce \unsure{propos}, je tombe sur le truc suivant, qui pour l'instant reste heuristique. Si $M\unsure{,|.}N$ sont \unsure{deux} faisceaux abéliens sur un topos $X$, et $\tau_{\unsure{\leq|<}2}\mathbb{R}\underline{\mathrm{Hom}}(M,N)=E(M,N)$ est le complexe ayant les invariants
\[
\left\{\begin{aligned}
\underline{H}^{\unsure{i|0}} &= \deleted{\unsure{\mathrm{Hom}}}\underline{\mathrm{Ext}}^{i}(M,N)  && \text{pour}\quad 0\leq i\leq 2\\
\underline{H}^{\unsure{i|j}} &= 0 && \text{si}\quad i\notin[0,2]
\end{aligned}\right.
\]
il doit y avoir un triangle distingué canonique
\diagram{\begin{tikzcd}[ampersand replacement=\&,column sep=small,row sep=small]
\& \& {\underline{\mathrm{Hom}}(M,{}_{2}N)[-2]}\arrow[dl] \& \\
(T) \& {E(M,N)}\arrow[rr] \& \& {E'(M,N)}\arrow[ul]
\end{tikzcd}}

% archive sheet number: 61

\scanpage{5}
% Two short horizontal strokes above the first line.
\margin{\mbox{-- --}}

donc $E'(M\;N)$ est un complexe dont les invariants $\underline{H}^{i}$ sont ceux de $E(M,N)$ en degré $i\unsure{+|\ne}2$, et qui en degré $2$ donne lieu \unsure{:|à} une suite exacte
\begingroup
\makeatletter\tagsleft@true\makeatother
\[
\tag*{(\unsure{*|S})}
0\to\underline{\mathrm{Ext}}^{2}(M,N)\to
\overbrace{\underline{H}^{2}\!\left(E'(M,N)\right)}^{P(M,N)}
\xrightarrow{\sigma}\underline{\mathrm{Hom}}(M,{}_{2}N)\to \unsure{0}
\]
\endgroup

Heuristiquement, $E'(M,N)$ est le complexe qui exprime le ``2-champ de Picard\deleted{\unsure{''|*}} strict'' formé des 1-champs de Picard (pas \unsure{nécessairement} stricts) \inserted{``épinglés'' par $M,N$} sur des objets variables de $X$, en admettant que ta théorie pour les 1-champs de Picard \inserted{stricts} s'étend aux 2-champs de Picard stricts (ce qui pour moi ne fait guère de doute) ; de même $E(M,N)$ correspond aux champs de Picard \uline{stricts} épinglés par $M,N$.

La suite exacte $(*)$ se construit en \unsure{tous|tout} cas canoniquement ``à la main'', où \deleted{\unsure{le}} terme médian est le faisceau des classes \inserted{à} \deleted{\unsure{d'|à}}``équivalence'' près des

% Two curved impressions run along the lower left edge of the copied sheet.
% archive sheet number: 62

\scanpage{6}
% Two short horizontal strokes at the upper right.
\margin{\mbox{-- --}}

\deleted{\unsure{catégories|catégorie}}\inserted{champs} de Picard \unsure{épinglés|épinglées} par $M\unsure{.|,}N$, $\sigma$ étant l'invariant qui \inserted{s'obtient en\unsure{,}} associant à toute section \inserted{$L$} d'un champ de Picard la symétrie de $\deleted{\unsure{N|L}}L\otimes L$, \unsure{interprétée} comme section de ${}_{2}N$. Je sais prouver (sauf erreur) que tout hom. $M\to{}_{2}N$ provient d'un champ de Picard convenable (épinglé par $M,N$) (a priori l'obstruction est dans $\mathrm{Ext}^{3}(X;M,N)$, mais un argument ``\unsure{universel|universal}'' prouve qu'elle est nulle). Cela prouve que l'extension $(*)$ est bien proche d'être splittée : \deleted{\unsure{la|un}} tte section du troisième faisceau, sur \unsure{un} objet quelconque de $X$, se \unsure{remonte} -- en d'autres termes l'extension a une section ``ensembliste''. Bien sûr, il y a mieux en fait : toute section sur un $U\in\mathrm{Ob}\,X$ ``provient'' d'un élément de ${\unsure{\mathscr{H}|\mathcal{H}}}^{2}(U,{\unsure{E|\underline{F}}}'(M\;N\unsure{,|)}))$ (hypercohomologie $\mathscr{H}^{2}$).

% archive sheet number: 63

\scanpage{7}

\uline{\illegible{1}}\unsure{:|,} \unsure{soit} $\unsure{\mathcal{A}|\mathscr{A}|A}$ un anneau sur $X$, soient $M,N$ respectivement les faisceaux $K^{0}$ et $K^{1}$ associés au champ additif des $\mathcal{A}$-Modules projectifs de \unsure{v.|t.|r.} f. \inserted{(\unsure{p|par} ex.),} Alors la \unsure{construction} de Mme Sinh nous \unsure{fourni|fournit|fournir} un champ de Picard épinglé par $M,N$, d'où une section \deleted{\unsure{du|de}} canonique du terme médian $\unsure{P|F}(M,N$ de $(*)$.

\uline{NB} Tout ce qui précède a les \unsure{fonctorialités} évidentes en $M,N,X$\ldots

\uline{Question} : Le triangle exact $(\unsure{T|J})$ \unsure{et} la suite exacte $(\unsure{*})$ sont ils \unsure{connu|connus} par les \unsure{compétences|compétents} (Quillen, Breen Illusie\ldots) ? \unsure{Connaissent-ils|Donnaissent-ils} des \unsure{variantes|variante} ``supérieures'' ? (Un principe ``géométrique'' pour en obtenir pourrait être via des $n$-champs de Picard non \unsure{néc.} stricts\ldots)

% Two curved impressions run along the lower left edge of the copied sheet.
% archive sheet number: 64

